% QUESTAO
Um aluno apresentou durante o ano letivo o seguinte aproveitamento: 1º Período: nota 6; 2º Período: nota 8; 3º Período: nota 6; e 4º Período: nota 9. Qual gráfico de linhas melhor representa essa situação?

% ALTERNATIVAS
\parbox{\linewidth}{\centering\tikz[scale=0.3]{\draw[CorSerie!30, xstep=3, ystep=1] (-0.5,-0.5) grid (9.5,3.5);\draw[ultra thick, CorSerie] (0,0)--(3,2)--(6,0)--(9,3);\draw[thick, CorSerie] (-0.5,-0.5) rectangle (9.5,3.5);\node[right, CorSerie] at (9.5,0) {\footnotesize 6};\node[right, CorSerie] at (9.5,3) {\footnotesize 9};}}
\parbox{\linewidth}{\centering\tikz[scale=0.3]{\draw[CorSerie!30, xstep=3, ystep=1] (-0.5,-0.5) grid (9.5,3.5);\draw[ultra thick, CorSerie] (0,0)--(3,3)--(6,1)--(9,1);\draw[thick, CorSerie] (-0.5,-0.5) rectangle (9.5,3.5);\node[right, CorSerie] at (9.5,0) {\footnotesize 6};\node[right, CorSerie] at (9.5,3) {\footnotesize 9};}}
\parbox{\linewidth}{\centering\tikz[scale=0.3]{\draw[CorSerie!30, xstep=3, ystep=1] (-0.5,-0.5) grid (9.5,3.5);\draw[ultra thick, CorSerie] (0,0)--(3,1)--(6,3)--(9,2);\draw[thick, CorSerie] (-0.5,-0.5) rectangle (9.5,3.5);\node[right, CorSerie] at (9.5,0) {\footnotesize 6};\node[right, CorSerie] at (9.5,3) {\footnotesize 9};}}
\parbox{\linewidth}{\centering\tikz[scale=0.3]{\draw[CorSerie!30, xstep=3, ystep=1] (-0.5,-0.5) grid (9.5,3.5);\draw[ultra thick, CorSerie] (0,3)--(3,0)--(6,0)--(9,1);\draw[thick, CorSerie] (-0.5,-0.5) rectangle (9.5,3.5);\node[right, CorSerie] at (9.5,0) {\footnotesize 6};\node[right, CorSerie] at (9.5,3) {\footnotesize 9};}}
\parbox{\linewidth}{\centering\tikz[scale=0.3]{\draw[CorSerie!30, xstep=3, ystep=1] (-0.5,-0.5) grid (9.5,3.5);\draw[ultra thick, CorSerie] (0,3)--(3,3)--(6,0)--(9,3);\draw[thick, CorSerie] (-0.5,-0.5) rectangle (9.5,3.5);\node[right, CorSerie] at (9.5,0) {\footnotesize 6};\node[right, CorSerie] at (9.5,3) {\footnotesize 9};}}


